% !TeX program = xelatex
% !TeX encoding = UTF-8
\documentclass[12pt,a4paper]{article}
\usepackage[margin=2.5cm]{geometry}
\usepackage{amsmath,amssymb,amsthm}
\usepackage{xeCJK}
\setCJKmainfont{Microsoft JhengHei}
\usepackage{graphicx}
\usepackage[hidelinks]{hyperref}
\usepackage{fancyhdr}
\usepackage{needspace}
\makeatletter
\newcommand{\totalpages}{\@abspage@last}
\makeatother
\pagestyle{fancy}
\fancyhf{}
\renewcommand{\headrulewidth}{0pt}
\fancyfoot[C]{\thepage/\totalpages}
\fancypagestyle{plain}{%
  \fancyhf{}%
  \renewcommand{\headrulewidth}{0pt}%
  \fancyfoot[C]{\thepage/\totalpages}%
}
\setlength{\parindent}{0pt}
\setlength{\parskip}{0.5em}
\raggedbottom
\title{Linear System --- Homework 4}
\author{葉彥辰\quad P46154345}
\date{}
\begin{document}
\maketitle
\begin{center}Fall 2026\quad Due: Friday, Oct 23\end{center}

\section*{Exercise 1}
Determine whether or not the following matrix is nondefective.
\[
A=\begin{pmatrix}0&1&0\\0&0&1\\-1&1&1\end{pmatrix}
\]
\subsection*{Solution}
% 在這裡撰寫解答。

\begin{align*}
	\det (A-sI) &= \begin{bmatrix}
		-s & 1 & 0 \\
		0 & -s & 1 \\
		-1 & 1 & 1 - s
	\end{bmatrix} \\
	&= -s \bigg( s(s-1) -1 \bigg ) + (-1)\cdot 1  \\
	&= -s(s^2 -s -1) - 1 \\
	&= -s^3 + s^2 + s - 1 \\
	&= -(s-1)^2(s+1) 
\end{align*}

for $s = 1$
\begin{align*}
	A - sI = \begin{bmatrix}
		-1 & 1 & 0 \\
		0 & -1 & 1 \\
		-1 & 1 & 0
	\end{bmatrix} \Rightarrow \mathbf v = \begin{bmatrix}
	1 \\ 1\\1
	\end{bmatrix}
\end{align*}
Hence, for eigenvalue $s = 1$, AM = 2, and GM = 1

 
for $s = -1$
\begin{align*}
	A - sI = \begin{bmatrix}
		1 & 1 & 0 \\
		0 & 1 & 1 \\
		-1 & 1 & 2\\
	\end{bmatrix} \Rightarrow \mathbf v = \begin{bmatrix}
	1 \\ -1 \\ 1
	\end{bmatrix}
\end{align*}

Hence, for eigenvalue $ s= -1$, AM = 1, and GM = 1

Therefore, matrix $A$ is defective.
\vspace{3cm}

\section*{Exercise 2}
What is the companion matrix whose eigenvalues are $-1$, $-2$, and $-3$?
\subsection*{Solution}
% 在這裡撰寫解答。

\begin{align*}
	p(s) &= (s + 1 )(s + 2) (s + 3) \\
	&= a_3 s^3 + a_2 s^2 + a_1s + a_0 \\
	&= s^3 + 6s^2 + 11 s + 6 \\
\end{align*}
Hence, companion matrix is 
\begin{align*}
	A &= \begin{bmatrix}
		0 & 1 & 0 \\
		0 & 0 & 1 \\
		-a_0 & -a_1 & -a_2
	\end{bmatrix} \\
	&= \begin{bmatrix}
		0 & 1 & 0 \\
		0 & 0 & 1 \\
		-6 & -11 & -6
	\end{bmatrix}
\end{align*}

\section*{Exercise 3}
What is the real $2\times2$ companion matrix with eigenvalues $1+j$, $1-j$?
\subsection*{Solution}
% 在這裡撰寫解答。
\begin{align*}
	p(s) &= (s - (1 + j)) (s - (1 - j)) \\
	&= a_2 s^2 + a_1 s + a_0 \\
	&=  s^2 - 2s + 2
\end{align*}
Hence, companion matrix is  
\begin{align*}
	A &= \begin{bmatrix}
		0 & 1 \\
		-a_0 & -a_1
	\end{bmatrix} \\
	&= \begin{bmatrix}
		0 & 1 \\
		-2 & 2
	\end{bmatrix}
\end{align*}

\vspace{3cm}

\section*{Exercise 4}
Suppose $A$ is a real square matrix and the vectors
\[
\begin{pmatrix}1\\-1\\1\\-1\end{pmatrix},\qquad
\begin{pmatrix}1\\2\\4\\8\end{pmatrix},\qquad
\begin{pmatrix}1\\j\\-1\\-j\end{pmatrix}
\]
are eigenvectors of $A$ corresponding to eigenvalues $-1$, $2$, and $j$, respectively. What is the response $x(t)$ of the system
\[\dot{x}=Ax\]
to the following initial conditions?

\textbf{(a)} $x(0)=\begin{pmatrix}2\\-2\\2\\-2\end{pmatrix}$
\subsection*{Solution}
% 在這裡撰寫解答。

Since $\mathbf v_1$ is an eigenvector corresponding to eigenvalue $\lambda_1=-1$,
\begin{align*}
	x(0) &= \begin{bmatrix}2\\-2\\2\\-2\end{bmatrix}
	= 2\begin{bmatrix}1\\-1\\1\\-1\end{bmatrix}
	= 2\mathbf v_1 \triangleq \mathbf v \\
\end{align*}
\[
A\mathbf v = \lambda \mathbf v \Leftrightarrow x(t)  = e^{\lambda t} \mathbf v= e^{\lambda t} x(0)
\]
\begin{align*}
	x(t) &= e^{\lambda_1 t} x(0) \\
	 &= e^{\lambda_1 t} \mathbf v \\
		&= e^{-t} \cdot 2 \mathbf v_1 \\
		&= \begin{bmatrix}2e^{-t}\\-2e^{-t}\\2e^{-t}\\-2e^{-t}\end{bmatrix}.
\end{align*}


\textbf{(b)} $x(0)=\begin{pmatrix}-1\\-2\\-4\\-8\end{pmatrix}$
\subsection*{Solution}
% 在這裡撰寫解答。

\[
x(0) = \begin{bmatrix}
	-1 \\ -2 \\ -4 \\ -8 
\end{bmatrix} = -\begin{bmatrix}
	1 \\ 2 \\ 4 \\ 8
\end{bmatrix} = -\mathbf v_2 \text{ which corresponding to eigenvalue is } 2
\]
\[
A \mathbf v = \lambda \mathbf v \Leftrightarrow x(t) = e^{\lambda t} x(0) = e^{\lambda t} \mathbf v
\]
\begin{align*}
	x(t) &= e^{\lambda t} x(0) \\
	&= e^{\lambda t} \mathbf v \\
	&= e^{\lambda t} \cdot - \mathbf v_2 \\
	&= e^{2t} \begin{bmatrix}
		-1 \\ -2 \\ -4 \\ -8 
	\end{bmatrix} = \begin{bmatrix}
	-e^{2t} \\ -2 e^{2t} \\ -4 e^{2t} \\ -8 e^{2t}
	\end{bmatrix}
\end{align*}

\textbf{(c)} $x(0)=\begin{pmatrix}0\\3\\3\\9\end{pmatrix}$
\subsection*{Solution}
% 在這裡撰寫解答。
\[
x(0) = \mathbf v = \begin{bmatrix}
	0 \\ 3 \\ 3 \\ 9
\end{bmatrix} = -\begin{bmatrix}
1 \\ -1 \\ 1 \\ -1
\end{bmatrix} + \begin{bmatrix}
1 \\ 2 \\ 4 \\8
\end{bmatrix} = -\mathbf v_1 + \mathbf v_2
\]
with linear combination
\begin{align*}
	x(t) &= -e^{\lambda_1 t}\mathbf v_1 + e^{\lambda_2 t} \mathbf v_2 \\
	&= \begin{bmatrix}
		-e^{-t} + e^{2t} \\ e^{-t}+2e^{2t} \\ -e^{-t }+4e^{2t} \\ e^{-t} + 8e^{2t}
	\end{bmatrix} 
\end{align*}

\textbf{(d)} $x(0)=\begin{pmatrix}1\\1\\-1\\-1\end{pmatrix}$
\subsection*{Solution}
% 在這裡撰寫解答。

\[
\mathbf v_3 = \begin{bmatrix}
	1 \\ j \\ -1 \\ j
\end{bmatrix} = \mathbf u_3 + j \mathbf p_3 = \begin{bmatrix}
1 \\ 0 \\ -1 \\ 0
\end{bmatrix} + j \begin{bmatrix}
0 \\ 1 \\ 0 \\ 1 
\end{bmatrix} \Rightarrow \mathbf u_3 = \begin{bmatrix}
1 \\ 0 \\ -1 \\ 0
\end{bmatrix}, \quad \mathbf p_3 = \begin{bmatrix}
0 \\ 1 \\ 0 \\ 1 
\end{bmatrix}
\]
\[
\lambda_3 = \alpha + j \omega = j \Rightarrow \alpha = 0, \quad \omega = 1
\]

\begin{align*}
	x(t) &= e^{\lambda t} \math
\end{align*}
\vspace{3cm}

\section*{Exercise 5}
Consider a discrete-time LTI system described by
\[x(k+1)=Ax(k).\]
\textbf{(a)} Suppose that the vectors
\[
\begin{pmatrix}1\\-1\\1\\-1\end{pmatrix}\qquad\text{and}\qquad
\begin{pmatrix}1\\1\\1\\1\end{pmatrix}
\]
are eigenvectors of $A$ corresponding to eigenvalues $-2$ and $3$, respectively. What is the response $x(k)$ of the system to the initial condition
\[x(0)=\begin{pmatrix}1\\0\\1\\0\end{pmatrix}?\]
\subsection*{Solution}
% 在這裡撰寫解答。
\vspace{3cm}

\textbf{(b)} Suppose $A$ is a real matrix and the vector
\[\begin{pmatrix}1\\j\\-1\\-j\end{pmatrix}\]
is an eigenvector of $A$ corresponding to the eigenvalue $2+3j$. What is the response $x(k)$ of the system to the initial condition
\[x(0)=\begin{pmatrix}1\\0\\-1\\0\end{pmatrix}?\]
\subsection*{Solution}
% 在這裡撰寫解答。
\vspace{3cm}

\section*{Exercise 6}
Consider a $5\times5$ matrix
\[A=vw^T\]
where $v$ and $w$ are $5\times1$ matrices.

\textbf{(a)} What is the rank of $A$?
\subsection*{Solution}
% 在這裡撰寫解答。
\vspace{3cm}

\textbf{(b)} What is the nullity of $A$?
\subsection*{Solution}
% 在這裡撰寫解答。
\vspace{3cm}

\textbf{(c)} What are the eigenvalues of $A$ if $w^Tv\neq0$?
\subsection*{Solution}
% 在這裡撰寫解答。
\vspace{3cm}

\clearpage
\textbf{For Exercise 7 and 8, please use Matlab/Python to illustrate the time history plot. Please don't plot $y$, $\theta_1$, $\theta_2$ in the same plot, separate them!}

\section*{Exercise 7}
Recall the 2 pendulum cart system.
\begin{align*}
&(m_0+m_1+m_2)\ddot{y}-m_1l\cos\theta_1\,\ddot{\theta}_1-m_2l\cos\theta_2\,\ddot{\theta}_2\\
&\qquad{}+m_1l\sin\theta_1\,\dot{\theta}_1^2+m_2l\sin\theta_2\,\dot{\theta}_2^2=u\\
&-m_1l\cos\theta_1\,\ddot{y}+m_1l^2\ddot{\theta}_1+m_1lg\sin\theta_1=0\\
&-m_2l\cos\theta_2\,\ddot{y}+m_2l^2\ddot{\theta}_2+m_2lg\sin\theta_2=0
\end{align*}
Consider the equilibrium configurations defined by $u^e=0$ and
\[E1:\qquad(y^e,\theta_1^e,\theta_2^e)=(0,0,0)\]
and
\[E2:\qquad(y^e,\theta_1^e,\theta_2^e)=(0,\pi,\pi).\]
Consider state space representations of the linearizations corresponding to the following combinations of parameters and equilibrium conditions:
\begin{center}
\begin{tabular}{c|c|c}
 & Parameter & Equilibrium\\\hline
$L1$ & $P1$ & $E1$\\
$L2$ & $P1$ & $E2$\\
$L3$ & $P2$ & $E1$\\
$L4$ & $P2$ & $E2$\\
$L5$ & $P3$ & $E1$\\
$L6$ & $P3$ & $E2$\\
$L7$ & $P4$ & $E1$\\
$L8$ & $P4$ & $E2$
\end{tabular}
\end{center}

\textbf{(a)} Determine the eigenvalues of all the \textbf{\textit{linearized systems}} L1--L8.
\subsection*{Solution}
% 在這裡撰寫解答。
\vspace{3cm}

\textbf{(b)} Compare the behavior of the nonlinear system with that of the linearized system for cases L7 and L8. Illustrate your results using different initial condition with time histories of $y$, $\theta_1$, and $\theta_2$. Plesae write paragraph to analysis the comparison between nonlinear and linearized system.

Initial conditions for case L7: $\pm0$ deg, $\pm1$ deg, $\pm10$ deg, $\pm90$ deg, $\pm180$ deg.

Initial conditions for case L8: $\pm1$ deg.
\subsection*{Solution}
% 在這裡撰寫解答。
\vspace{3cm}

\section*{Exercise 8}
This exercise refers to linearizations L7 and L8 of the two pendulum cart system.

\textbf{(a)} For L7 choose an initial state for the linearized system which results in a periodic solution for the linearized system.
\subsection*{Solution}
% 在這裡撰寫解答。
\vspace{3cm}

\textbf{(b)} For L8 choose an initial state for the linearized system which results in a solution which asymptotically decays to zero for the linearized system.
\subsection*{Solution}
% 在這裡撰寫解答。
\vspace{3cm}

\textbf{(c)} For L8 choose an initial state for the linearized system which results in a solution whose magnitude grows exponentially for the linearized system.
\subsection*{Solution}
% 在這裡撰寫解答。
\vspace{3cm}

In each case, simulate both the linearized system and the nonlinear system with initial conditions corresponding to your chosen initial state for the linearized system.
\end{document}
